\chapter{StochasticSimulator}
\label{chap:stochastic_simulator}

\section{Purpose}

The \texttt{StochasticSimulator} class is the core simulation engine of
\actsim. It supports Monte Carlo simulation of aggregate losses under
specified frequency and severity distributions and can also be used to
generate correlated multivariate simulations across multiple lines of
business.

The simulator is designed to be reproducible: every run is governed by an
explicit random seed, and intermediate event-level data can be retained
for downstream metrics such as Occurrence Exceedance Probability.

\section{Construction}

A simulator is constructed by specifying the frequency and severity
distributions together with their parameters:

\begin{lstlisting}[style=pythonstyle]
from actsim import StochasticSimulator

simulator = StochasticSimulator(
    freq_dist="poisson",
    freq_params=(10,),
    sev_dist="lognormal",
    sev_params=(10, 0.5),
    num_sim=10000,
    keep_all=True,
    seed=1234,
)
\end{lstlisting}

The constructor arguments are summarised below.

\begin{center}
\begin{tabular}{ll}
\toprule
\textbf{Argument} & \textbf{Description} \\
\midrule
\texttt{freq\_dist}    & Name of the frequency distribution. \\
\texttt{freq\_params}  & Tuple of parameters for the frequency distribution. \\
\texttt{sev\_dist}     & Name of the severity distribution. \\
\texttt{sev\_params}   & Tuple of parameters for the severity distribution. \\
\texttt{num\_sim}      & Number of simulation periods (default 10{,}000). \\
\texttt{keep\_all}     & If \texttt{True}, retain individual event records. \\
\texttt{seed}          & Random seed for reproducibility. \\
\texttt{correlation}   & Optional linear correlation between $N$ and $X$. \\
\texttt{copula\_type}  & Optional copula family. \\
\texttt{theta}         & Copula dependence parameter. \\
\bottomrule
\end{tabular}
\end{center}

\section{Frequency--Severity Simulation}

The primary method is \texttt{gen\_agg\_simulations()}, which loops over
the requested number of simulation periods, draws a frequency from the
frequency distribution, and then draws that many severities from the
severity distribution. Aggregate losses are accumulated and stored in
\texttt{results}.

\begin{lstlisting}[style=pythonstyle]
simulator.gen_agg_simulations()
print(simulator.results.head())
\end{lstlisting}

If \texttt{keep\_all=True}, individual event-level records are stored in
\texttt{all\_simulations}, with columns \texttt{year}, \texttt{event\_id},
\texttt{yearly\_event\_id}, and \texttt{amount}.

\section{Dependency between Frequency and Severity}

\subsection{Linear Correlation}

When \texttt{correlation} is supplied without a copula, frequency and
severity are correlated through a Cholesky factorisation of a $2 \times 2$
correlation matrix:
\[
    C = \begin{pmatrix} 1 & \rho \\ \rho & 1 \end{pmatrix},
    \qquad C = L L^{\top}.
\]
Standard normal variates are transformed into uniform variates via the
normal CDF, and these uniforms are then mapped through the inverse CDF of
the target frequency and severity distributions.

\begin{lstlisting}[style=pythonstyle]
simulator = StochasticSimulator(
    "poisson", (10,),
    "lognormal", (10, 0.5),
    num_sim=10000,
    keep_all=True,
    seed=1234,
    correlation=0.6,
)
\end{lstlisting}

\subsection{Copula-Based Dependency}

\actsim supports four copula families: Gaussian, Frank, Gumbel, and
Clayton. Each is parametrised by a dependence parameter
$\theta$ (or by a correlation matrix in the Gaussian case).

\begin{lstlisting}[style=pythonstyle]
simulator = StochasticSimulator(
    "poisson", (10,),
    "lognormal", (10, 0.5),
    num_sim=10000,
    keep_all=True,
    seed=1234,
    correlation=0.6,
    copula_type="frank",
    theta=0.6,
)
\end{lstlisting}

The chosen copula generates pairs of dependent uniforms, which are then
mapped to frequency and severity via the inverse CDFs.

\section{Multivariate Correlated Simulation}

The method \texttt{gen\_multivariate\_corr\_simulations} generates
correlated losses across multiple lines of business. It accepts a
correlation matrix as a CSV file and a list of marginal distributions as
a JSON file.

\begin{lstlisting}[style=pythonstyle]
simulator = StochasticSimulator(
    "normal", (1, 0),
    "normal", (1, 0),
    num_sim=100000,
    keep_all=True,
    seed=1234,
)

simulator.gen_multivariate_corr_simulations(
    "examples/correlated_sim/corr_matrix.csv",
    "examples/correlated_sim/dist_list.json",
    gen_marginal=True,
)
\end{lstlisting}

When \texttt{gen\_marginal=True}, the simulator stores marginal
simulations for each line of business in \texttt{\_all\_simulations\_data}.
The aggregated total across all lines is stored in \texttt{results}.

\section{Analysing Results}

\subsection{Aggregate Percentiles}

The \texttt{calc\_agg\_percentile} method returns a single percentile of
the aggregate loss distribution.

\begin{lstlisting}[style=pythonstyle]
p99_2 = simulator.calc_agg_percentile(99.2)
print(f"99.2% percentile: {p99_2:,.2f}")
\end{lstlisting}

\subsection{Risk Metrics}

The \texttt{analyze\_results} method computes a table of VaR, TVaR, OEP,
and AEP at user-specified quantiles. OEP is only available when
\texttt{keep\_all=True}.

\begin{lstlisting}[style=pythonstyle]
metrics = simulator.analyze_results(
    quantiles=[0.5, 0.75, 0.9, 0.95, 0.99]
)
print(metrics)
\end{lstlisting}

\subsection{Distribution Plot}

The \texttt{plot\_distribution} method draws a histogram of the simulated
aggregate losses, optionally on a logarithmic scale:

\begin{lstlisting}[style=pythonstyle]
simulator.plot_distribution(bins=100, log_option=False)
\end{lstlisting}

\subsection{Correlation Plot}

If \texttt{keep\_all=True}, the \texttt{plot\_correlated\_variables}
method draws a scatter and density plot of yearly event counts against
mean severity, displaying the empirical Pearson correlation. This plot
is useful for sanity-checking whether the imposed dependency structure
behaves as expected.

\begin{lstlisting}[style=pythonstyle]
simulator.plot_correlated_variables()
\end{lstlisting}

\section{Applying Deductibles and Limits}
\label{sec:deductibles_and_limits}

Insurance contracts often include per-occurrence and aggregate
deductibles and limits. These can be applied to simulated losses with
\texttt{apply\_deductible\_and\_limit}:

\begin{lstlisting}[style=pythonstyle]
gross_loss = simulator.apply_deductible_and_limit(
    per_occurrence_ded=1000,
    per_occurrence_limit=10000,
    agg_ded=100000,
    agg_limit=300000,
)
\end{lstlisting}

The method first reduces each event's loss by the per-occurrence
deductible (with a floor at zero), caps the result at the per-occurrence
limit, then aggregates by year and applies the annual aggregate
deductible and limit.

The returned DataFrame contains the gross (capped and floored) annual
losses. To run further analyses on the post-contract losses, supply the
DataFrame as the \texttt{all\_simulations} keyword argument to
\texttt{analyze\_results}:

\begin{lstlisting}[style=pythonstyle]
gross_loss["amount"] = gross_loss["gross_loss"]
metrics = simulator.analyze_results(all_simulations=gross_loss)
\end{lstlisting}

\section{Reproducibility}

All random draws inside \texttt{StochasticSimulator} are governed by the
\texttt{seed} argument supplied at construction. Setting the seed to a
fixed value guarantees that the same input produces the same output
across runs, which is essential for validation and audit. The
recommended practice is to record the seed alongside any results that
are reported externally.
